\documentclass[../BTL.tex]{subfiles}
\begin{document}

\chapter{TỔNG QUAN VỀ ĐỀ TÀI}

\section{Bối cảnh lý do chọn đề tài}
Trong bối cảnh xã hội hiện đại, nhu cầu chăm sóc ngoại hình (Grooming)
của nam giới đang gia tăng mạnh mẽ, đặc biệt là tại các đô thị lớn. Khái
niệm "làm đẹp" giờ đây không chỉ dành cho phái nữ, mà phái mạnh cũng
ngày càng chú trọng đến phong cách tóc và râu. Tuy nhiên, với nhịp sống
bận rộn, quỹ thời gian của khách hàng ngày càng eo hẹp. Phương thức vận
hành truyền thống tại các \textbf{Barbershop} -- nơi khách hàng phải đến
trực tiếp ngồi chờ đến lượt ("first come, first served") -- đã bộc lộ
nhiều hạn chế, gây lãng phí thời gian và giảm trải nghiệm của những quý
ông bận rộn.

Mặt khác, đối với các chủ Barbershop, việc quản lý lịch hẹn thủ công
hoặc qua tin nhắn rời rạc thường dẫn đến tình trạng dồn khách vào giờ
cao điểm, trong khi giờ thấp điểm lại vắng khách, gây khó khăn trong
việc tối ưu hóa năng suất của các thợ cạo (Barber).

Đề tài \textbf{``Xây dựng Website đặt lịch cắt tóc cho Barbershop''}
được thực hiện nhằm giải quyết bài toán trên: giúp khách hàng chủ động
thời gian thông qua việc đặt chỗ trực tuyến (Booking), đồng thời cung
cấp cho Barbershop công cụ quản lý chuyên nghiệp. Đây là bước tiến cần
thiết để nâng cao chất lượng dịch vụ và định hình phong cách phục vụ
hiện đại cho phái mạnh.

\section{Mục đích của website cần thiết kế}
Đề tài "Xây dựng website đặt lịch cắt tóc" được thực hiện nhằm các mục
tiêu sau:

\begin{itemize}
\item
  Xây dựng hệ thống trực tuyến hỗ trợ khách hàng tự do chọn dịch vụ,
  chọn thợ và khung giờ phù hợp.
\item
  Đảm bảo trải nghiệm người dùng thân thiện, giao diện mượt mà không cần
  tải lại trang (mô hình SPA).
\item
  Tạo giao diện quản trị dành cho Admin để quản lý toàn diện lịch hẹn,
  khách hàng và đội ngũ thợ cắt tóc (Stylist).
\item
  Cung cấp giao diện riêng biệt cho thợ cắt tóc để theo dõi lịch làm
  việc cá nhân.
\item
  Ứng dụng kiến trúc RESTful API, chia tách rõ ràng Frontend (HTML/JS)
  và Backend (PHP) để dễ dàng bảo trì và mở rộng.
\end{itemize}

\section{Phạm vi nghiên cứu}
Đề tài tập trung phát triển hệ thống đặt lịch và quản lý dành cho ba đối
tượng chính: Khách hàng (Customer), Thợ cắt tóc (Stylist) và Quản trị
viên (Admin). Các chức năng chính bao gồm:

\begin{itemize}
\item
  \textbf{Cho Khách hàng:}

  \begin{itemize}
  \item
    Đăng ký, đăng nhập, đăng xuất tài khoản.
  \item
    Xem thông tin hồ sơ cá nhân.
  \item
    Xem danh sách chi nhánh, dịch vụ và thợ cắt tóc.
  \item
    Đặt lịch hẹn (hệ thống tự động tính toán và hiển thị giờ trống của
    thợ).
  \item
    Xem lịch sử đặt hẹn.
  \item
    Hủy lịch hẹn (khi lịch chưa được quản trị viên xác nhận).
  \end{itemize}
\item
  \textbf{Cho Thợ cắt tóc (Stylist):}

  \begin{itemize}
  \item
    Đăng nhập tài khoản nhân viên.
  \item
    Xem lịch làm việc (danh sách khách hàng đã đặt với mình theo
    ngày/giờ).
  \item
    Nộp đơn xin nghỉ phép.
  \end{itemize}
\item
  \textbf{Cho Quản trị viên (Admin):}

  \begin{itemize}
  \item
    Quản lý Chi nhánh: Thêm, sửa, xóa, xem danh sách các chi nhánh.
  \item
    Quản lý Dịch vụ: Thêm, sửa, xóa, xem danh sách dịch vụ cắt tóc.
  \item
    Quản lý Khách hàng: Thêm, sửa, xóa, tìm kiếm thông tin khách hàng.
  \item
    Quản lý Thợ cắt tóc: Thêm thợ mới, sửa thông tin, xóa/ngưng hoạt
    động tài khoản thợ.
  \item
    Quản lý Lịch hẹn: Xem danh sách toàn bộ lịch hẹn, cập nhật trạng
    thái (Chờ duyệt, Đã xác nhận, Đã hoàn thành, Đã hủy).
  \item
    Quản lý Nghỉ phép: Xem danh sách và xét duyệt đơn xin nghỉ phép của
    thợ.
  \end{itemize}
\end{itemize}

\section{Tổng kết chương}
Chương 1 đã trình bày bối cảnh, mục tiêu và phạm vi của đề tài, đồng
thời xác định rõ các đối tượng sử dụng và chức năng trọng tâm của hệ
thống. Đây là nền tảng để chuyển sang chương sau, nơi mô tả phương pháp
thực hiện, công nghệ sử dụng và công cụ hỗ trợ triển khai.
\end{document}
